\appendix

\section{对拟合恒星参数的敏感性 Sensitivity to Fitting the Stellar Parameters}
\label{app:stellar_label_sensitivity}

主 APOGEE 实验保留了 DR14 发表的、经光度定标的 $T_{\rm eff}$ 与经星震定标的巨星 $\log g$，并固定 $\xi$，以单独考察化学与仪器拟合 \citep{Holtzman2018}。我们用相同的官方光谱、样本、谱线定标、仪器算子、掩模、不确定度与连续谱处理来检验该选择：从主丰度解出发，把三个恒星参数与同样的 12 个丰度坐标一起拟合，并收敛最终的物理大气。图~\ref{fig:appendix_stellar_label_sensitivity} 与图~\ref{fig:appendix_abundance_sensitivity} 把所得恒星与化学分布同固定参数解及 APOGEE DR14 比较。

\begin{EN}
The main APOGEE experiment retains the photometrically calibrated $T_{\rm eff}$
and asteroseismically calibrated giant $\log g$ published in DR14, and fixes
$\xi$ to isolate the chemical and instrumental fit \citep{Holtzman2018}. We
test that choice with the same official spectra, sample, line calibration,
instrument operator, masks, uncertainties, and continuum treatment. Starting
from the main abundance solution, we fit those three stellar parameters
together with the same 12 abundance coordinates and converge the final physical
atmosphere. Figures
\ref{fig:appendix_stellar_label_sensitivity} and
\ref{fig:appendix_abundance_sensitivity} compare the resulting stellar and
chemical distributions with the fixed-parameter solution and APOGEE DR14.
\end{EN}

完整物理重拟合给出的逆选择权重中位变化为：$T_{\rm eff}$ $-200$ K、$\log g$ $+0.17$ dex、$\xi$ $-0.32$ km s$^{-1}$、$[\mathrm{Fe}/\mathrm{H}]$ $-0.06$ dex。在初始红团簇聚集区的 694 颗星中，中位变化为 $-220$ K 与 $+0.17$ dex。巨星支脊线与红团簇聚集仍清晰可见，但其绝对位置随所采用的恒星参数标度移动。

\begin{EN}
The full physical refit gives inverse-selection-weighted median changes of
$-200$~K in $T_{\rm eff}$,
$+0.17$~dex in $\log g$, $-0.32$~km~s$^{-1}$ in $\xi$, and $-0.06$~dex in
$[\mathrm{Fe}/\mathrm{H}]$. Among the 694 stars in the initial red-clump
concentration, the median changes are $-220$~K and $+0.17$~dex. The giant-branch
ridge and red-clump concentration remain visible, but their absolute location
moves with the adopted stellar-parameter scale.
\end{EN}

丰度中位变化从氮的 $-0.23$ dex、氧的 $-0.22$ dex 到硫的 $+0.22$ dex 不等；镁变化 $-0.03$ dex。丰度序列保持存在，而其零点响应恒星参数的变化。因此，更冷的光谱温度标度会传播为随元素变化的丰度零点。

\begin{EN}
The median abundance changes range from $-0.23$~dex for nitrogen and
$-0.22$~dex for oxygen to $+0.22$~dex for sulfur. Magnesium changes by
$-0.03$~dex. The abundance sequences persist while their zero points respond
to the stellar parameters. The cooler spectroscopic temperature scale therefore
propagates into element-dependent abundance zero points.
\end{EN}

生产级分析可以对光谱、宽波段能谱分布与视差做联合拟合，锚定 $T_{\rm eff}$ 与 $\log g$ 并吸收这一标度差异 \citep{Cargile2020}。本概念验证不包含这些外部信息。

\begin{EN}
A joint fit to the spectrum, broadband spectral energy distribution, and parallax can
anchor $T_{\rm eff}$ and $\log g$ and absorb this scale difference in a
production analysis \citep{Cargile2020}. We do not include that external
information in this proof of concept.
\end{EN}

\begin{physbox}{为什么恒星参数与丰度会“耦合”}
附录展示的是光谱分析中著名的\textbf{参数耦合}（或退化）现象：降低 $T_{\rm eff}$ 会让模型光谱整体变化，拟合器会调整丰度去补偿——例如温度影响电离平衡，从而改变电离与中性谱线的相对强度，进而影响推出的元素丰度。本附录定量说明：采用更冷的自由光谱温度（中位 $-200$ K）后，丰度零点整体移动最多约 0.2 dex，但族群的相对形态（双 $\alpha$ 序列、红团簇结构）不变。这说明本文主文结论对“固定 DR14 恒星参数”这一选择是稳健的，但绝对丰度标度需要外部信息（光度、视差）来锚定。
\end{physbox}

\begin{figure}[t]
\centering
\includegraphics[width=0.98\textwidth]{figure_14.pdf}
\captionof{figure}{$T_{\rm eff}$、$\log g$、$\xi$ 与 12 个丰度坐标联合拟合的敏感性。左：合并样本的逆选择权重等值线，灰虚线保留 DR14 恒星参数，深色实线为最终物理收敛拟合。右：固定参数的铁丰度为灰色空心圆，最终值按 APOGEE 高/低 $\alpha$ 样本着色。（Sensitivity to fitting the stellar parameters.）}
\label{fig:appendix_stellar_label_sensitivity}
\vspace{8pt}
\includegraphics[width=0.98\textwidth]{figure_15.pdf}
\captionof{figure}{$T_{\rm eff}$、$\log g$、$\xi$ 与丰度最终物理收敛拟合的逐元素丰度分布。布局与标度同图~\ref{fig:apogee_elements}；彩色点与等值线为两个 APOGEE 样本的最终值，灰虚等值线为 APOGEE DR14 的 50、80\% 密度水平。（Elemental-abundance distributions from the final physically converged fit.）}
\label{fig:appendix_abundance_sensitivity}
\end{figure}

